\documentclass[11pt]{article}

\usepackage[margin=1in]{geometry}
\usepackage{lmodern}
\usepackage[T1]{fontenc}
\usepackage[utf8]{inputenc}
\usepackage{amsmath, amssymb}
\usepackage{fancyvrb}
\usepackage{xcolor}
\usepackage{hyperref}
\hypersetup{colorlinks=true, urlcolor=blue, linkcolor=black}

\title{\textbf{Never Quantum at All}\\[6pt]
 \large Do we really need quantum computing?}
\author{Dmitry Sierikov}
\date{}

\newcommand{\Phione}{$\Phi_1$}

\begin{document}
\maketitle

\begin{abstract}
\noindent A student dequantized a famous quantum algorithm. A single matrix-free
dial---the effective rank---tells you which speedups survive a classical attack,
and which were never quantum at all. Every number below is reproduced by a
runnable demo in \texttt{stands/}, checked against ground truth.
\end{abstract}

\section*{The speedup that wasn't}

In 2018 a then-undergraduate, Ewin Tang, did something that was supposed to be
impossible. There was a celebrated quantum algorithm for recommendation
systems---the kind that suggests your next film---and it ran exponentially faster
than any known classical method. It was a poster child for ``quantum advantage in
machine learning.''

Tang \emph{dequantized} it. She wrote a \emph{classical} algorithm that matched the
quantum one's speed. The exponential gap vanished. And it was not a one-off: a
whole family of ``quantum machine learning'' speedups quietly collapsed the same
way.

So the uncomfortable question is real. When someone shows you a quantum speedup,
how do you know it is genuine---that no clever student will dequantize it next
year---versus a speedup that was \emph{redundant all along}, classical advantage
wearing a quantum costume? There is a single number you can measure, matrix-free,
that---in every case we have tested---separates the speedups with a known classical
attack from those that resist one. Not a proof; a diagnostic that has been right
wherever we have checked it.

\section*{One number: effective rank}

First, the one tool this piece leans on. \texttt{resona} is a small open-source
library for \emph{matrix-free} spectral computation: it answers questions about a
giant operator---its spectrum, its rank, its structure---using nothing but
matrix--vector products, never building or storing the matrix itself. That
restriction is the whole point here, because the operators in question are far too
big to write down. You hand \texttt{resona} a way to multiply by your operator; it
hands back the numbers.

The number we want from it is \textbf{effective rank}---in \texttt{resona},
\Phione. The shorthand: \emph{quantum advantage needs structure a classical machine
cannot cheaply simulate.} And the most common ``structure'' in quantum ML---low
rank---is exactly the structure classical sampling harvests for free.

The intuition: if a giant operator secretly lives in a low-dimensional subspace,
you do not need to touch the whole thing. You can \emph{sample} a few rows or
columns and reconstruct that subspace; cost scales with the rank, not the
dimension. That is precisely what Tang's algorithm does, and why the quantum
speedup evaporated: the quantum computer was paying (in entanglement, in qubits)
for structure that classical sampling gets at a discount.

\Phione{} measures that structure as a single scalar---the effective number of
directions the operator really uses.
\[
  \Phi_1 \;=\; \frac{\big(\mathrm{Tr}\,A\big)^2}{\mathrm{Tr}\,A^2}.
\]
Empirically the reading splits two ways: \textbf{low \Phione{} tracks
dequantizability}---a few directions carry the action, so sampling reconstructs
it; \textbf{high \Phione{} tracks the genuine frontier}---no low-rank handle for
sampling to grab. This is a measured correlation across the cases below, not a
theorem. A low reading flags a problem as a \emph{candidate} for classical attack;
a high reading says only that \emph{this probe} found no extractable
structure---not that none exists. One dial, computed from matrix--vector products,
with no matrix ever formed.

The reframing is worth pausing on. The usual question---``is this algorithm quantum
or classical?''---is the wrong one. The right question is ``how much
\emph{structure} does the problem have?'', and structure is measurable.

\begin{quote}
\textbf{A quantum computer pays a fixed price---qubits, coherence, error
correction---to deal with problems whose structure cannot be harvested
classically.}
\end{quote}

\noindent If the structure is low, you overpaid: a classical sampler does the same
job for free.

\section*{Where the advantage evaporates}

The stand is \texttt{stands/dequantize.py}. It builds an implicit low-rank operator
(true rank $r=5$, embedded in up to a million rows) and recovers its top singular
subspace by sampling a fixed $s=60$ rows:

\begin{Verbatim}[fontsize=\small]
   n (rows)  rows sampled   data touched   subspace overlap
 --------------------------------------------------------------
      2,000            60        3.0000%             1.0000
     20,000            60        0.3000%             1.0000
    200,000            60        0.0300%             1.0000
  1,000,000            60        0.0060%             1.0000
\end{Verbatim}

The same 60 samples work as $n$ grows $500\times$. At a million rows we touch
\textbf{0.006\% of the data} and still recover the subspace with overlap
\textbf{1.0000}. And the dial reads the situation correctly:

\begin{Verbatim}[fontsize=\small]
Phi_1 = resona.effective_rank(A^T A) = 3.28   (true rank r=5)  -- LOW => dequantizable
\end{Verbatim}

Low \Phione. The ``exponential quantum advantage'' here was redundant---it was
paying for low-rank structure that classical sampling extracts for free. One
fair-play note: ``data touched'' counts the sketched rows, and the
$\lVert\text{row}\rVert^2$ sampling distribution is assumed given---the
\emph{same} prepared-access assumption the quantum algorithm itself made (its
QRAM; classically, Tang's data structure maintains those norms as the data is
written).

\section*{Shor: where the wall is real}

Now point the \emph{same dial} at Shor's algorithm---factoring, the speedup that
actually threatens RSA. The stand is \texttt{stands/shor\_wall.py}. For small order
$r$, even classical period-finding factors $N$ cheaply; the trouble is that for
RSA-sized $N$ the order is itself exponentially large. So look at what \Phione{}
reads on the sequence $a^x \bmod N$---the thing Shor's quantum period-finding chews
on:

\begin{Verbatim}[fontsize=\small]
                    sequence   Phi_1 (eff. rank, max=60)
     a^x mod N (Shor target)                 19.4   <- ~10x higher: no handle
         periodic (period 7)                  2.0   <- low: we harvest it
\end{Verbatim}

The Shor sequence reads \textbf{$\Phi_1 = 19.4$}---about \textbf{$10\times$ higher}
than a simple periodic signal: no low-rank handle \emph{that this probe can see}.
And the extractability test---does the structure
\emph{saturate} as you widen the window (removable) or keep \emph{growing} (a
genuine wall)?---is decisive:

\begin{Verbatim}[fontsize=\small]
                    window ->    20     40     80    120   extractable?
     a^x mod N (Shor target)   1.5    6.1   24.7   40.3   False  <- GROWS: never closes
         periodic (period 7)   2.0    2.0    2.0    2.0   True   <- SATURATES
\end{Verbatim}

The periodic signal saturates at $2.0$---a finite chart closes, you harvest it
classically. The Shor target keeps growing ($1.5 \to 40.3$) and never closes: the
structure is \textbf{not extractable}, no finite window will ever contain it. That
is the stand's own verdict---the \texttt{extractable?\ False} in the table above is
its printed output. The same dial that flagged the low-rank problem
``redundant---sample it'' reads ``no extractable structure here'' for
Shor---\emph{from the same matrix-free measurement}. The first is a green light for
a classical attack; the second is the \emph{absence of a handle}, which is weaker
than a proof that no handle exists---but it is exactly where every genuine quantum
advantage we know of lives.

\section*{The honest limit}

One honest word on what \Phione{} is and isn't. It is not a proof that a speedup is
impossible---or that one is necessary. It measures a single thing: how much
low-dimensional structure a problem exposes to matrix-free probes. And the pattern
is empirical---worth saying plainly. Every dequantization we tested lands in the
low-\Phione{} regime; every textbook quantum advantage lands at the far end. That
makes \Phione{} a \emph{diagnostic} you can lean on---a number that has told the two
apart every time we checked---not a theorem that says they always must. A clean
track record, and the honesty to say ``so far.''

\Phione{} does not beat \emph{all} quantum computing---and the stand says so plainly.
It beats exactly the \textbf{low-rank / structured class}: low-rank linear algebra
(Tang), Clifford circuits, free fermions, low-entanglement dynamics, where
speedups are at most polynomial. It does \textbf{not} touch the genuine frontier:
Shor factoring and discrete log, real-time volume-law entanglement growth, generic
random-circuit sampling. Reading the dial as high does \emph{not} give you a
classical algorithm for Shor---beating Shor classically would mean factoring in
polynomial time and breaking RSA, and there is no known path. The framework
honestly points \emph{away} from that. The dial tells you which fight you are in;
it does not win the unwinnable one.

\section*{So --- do we need quantum computing?}

Yes---for the problems that earn it. Shor-style period-finding, volume-law
dynamics, generic random-circuit sampling: there the dial reads high, no
classical handle exists, and the quantum computer's price buys something real.
But for a surprising share of what has been sold as ``quantum advantage''---
everything living in the low-\Phione{} regime---the honest answer is the title
of this piece. Measure first. If the effective rank is low, you never needed a
quantum computer at all.

\section*{Run it yourself}

Everything---this article, both stands, and the figures---lives in one repository:
\url{https://github.com/dimaq12/do-we-need-quantum-computing}.

\begin{Verbatim}[fontsize=\small]
git clone https://github.com/dimaq12/do-we-need-quantum-computing.git
cd do-we-need-quantum-computing
pip install -r requirements.txt
python stands/dequantize.py
python stands/shor_wall.py
\end{Verbatim}

Every number above is reproducible: subspace overlap $1.0000$ at a million rows
touching $0.006\%$ of the data, $\Phi_1 = 3.28$ on the low-rank side, $\Phi_1 =
19.4$ and \texttt{extractable? False} on the Shor side. Don't trust me---run it.

\vspace{1em}
\noindent\rule{\linewidth}{0.4pt}\\
\footnotesize Built on \texttt{resona} (\url{https://pypi.org/project/resona/}),
matrix-free spectral computation. Part of the \emph{Spectra Without Matrices}
series.

\end{document}
